% Literature review for the MSc thesis "Where the Shortcut Sits" (requested by the supervisor, October 2026).
% Build: pdflatex literature_review && bibtex literature_review && pdflatex literature_review && pdflatex literature_review
% Bibliography: refs.bib (every 2021+ entry checked against its publisher or proceedings page; see Sec. 2).
\documentclass[11pt,a4paper]{article}
\usepackage[margin=2cm]{geometry}
\usepackage[T1]{fontenc}
\usepackage{lmodern}
\usepackage{amsmath,amssymb,booktabs,longtable,array,xcolor,enumitem}
\usepackage[round,authoryear]{natbib}
\usepackage[hidelinks]{hyperref}
\setlength{\emergencystretch}{2em}
\setlength{\parskip}{3pt}
\newcolumntype{P}[1]{>{\raggedright\arraybackslash}p{#1}}
\newcommand{\takeaway}[1]{\par\smallskip\noindent\colorbox{gray!12}{\parbox{\dimexpr\linewidth-2\fboxsep}{\small\textbf{For the thesis.} #1}}\par\medskip}
\newcommand{\bg}{\textsuperscript{\dag}}

\title{Shortcut learning from image artifacts, region-of-interest masking and its alternatives in medical image
classification: a review of the 2021--2026 literature}
\author{Literature review for the MSc thesis \emph{Where the Shortcut Sits}}
\date{October 2026}

\begin{document}
\maketitle

\begin{abstract}
\noindent This review collects the peer-reviewed work published from January 2021 to October 2026 that bears on one
question: when an image artifact that is correlated with the diagnosis lies inside or outside the region of interest
(ROI), does restricting the classifier to the ROI remove the shortcut, and what should be used instead? We included
68 studies from journals and conferences of recognised standing (Nature, Lancet and Science families, Medical Image
Analysis, MICCAI, CVPR, ECCV, NeurIPS, ICML, ICLR and others), verified each against its publisher or proceedings page
and excluded preprints and workshop papers. The literature establishes firmly that medical classifiers exploit
acquisition signals, devices, markings and demographic correlates; it documents artifacts per modality, mostly in
dermoscopy and chest radiography; it shows that image backgrounds carry predictive signal and that ROI restriction
sometimes helps; and it offers group-robustness, invariance, erasure, explanation-guided and generative mitigations,
none of which is uniformly better than standard training. Every masking and background-removal study we found assumes
that the shortcut lies outside the ROI. None treats the artifact's location relative to the ROI as an experimental
variable. Only one study, in natural images, reports that removing one shortcut can amplify another. The review ends
with a gap analysis that positions the thesis against the closest prior work and states what is and is not new.
\end{abstract}

\tableofcontents

%======================================================================================================================
\section{Purpose and review questions}\label{sec:purpose}
Deep classifiers for medical images often rely on features that predict the label in the training data without being
evidence of disease---\emph{shortcuts} \citep{geirhos2020shortcut}\bg. A common defence is to segment the lesion or
organ and classify only the region of interest. The thesis asks when that defence works, and finds that it depends on
where the artifact sits: masking removes an artifact outside the ROI but keeps, and makes the model rely more on, an
artifact inside it. This review sets out what was known before and alongside that work. It is organised by six
questions:
\begin{enumerate}[label=\textbf{RQ\arabic*},leftmargin=3em,itemsep=1pt]
\item What evidence is there that medical image classifiers use content that is not pathology (artifacts, devices,
acquisition signals, demographic correlates)? (Sec.~\ref{sec:evidence})
\item Which artifacts are documented in each modality, and where do they lie relative to the ROI? (Sec.~\ref{sec:modality})
\item What is known about restricting a model to the ROI by masking, segmentation or background removal? (Sec.~\ref{sec:masking})
\item How are shortcuts detected and audited? (Sec.~\ref{sec:detect})
\item Which mitigation methods exist, and what does each require? (Sec.~\ref{sec:mitigate})
\item What theory explains which feature a model uses, and how should claims of this kind be evaluated? (Secs.~\ref{sec:theory}, \ref{sec:eval})
\end{enumerate}
Section~\ref{sec:gap} draws the gap analysis and Section~\ref{sec:table} lists every included study.

%======================================================================================================================
\section{Method}\label{sec:method}
\paragraph{Search} Searches were run in September and October 2026 with web and scholarly search engines, PubMed, the
publishers' sites and the proceedings of MICCAI (Springer LNCS), CVF open access (CVPR, ICCV), ECVA (ECCV), PMLR (ICML,
AISTATS, MIDL), NeurIPS and OpenReview (ICLR). Query terms combined \emph{shortcut learning, spurious correlation,
Clever Hans, confounding, dataset bias} with \emph{medical imaging, dermoscopy, skin lesion, ultrasound, calipers,
chest radiograph, chest drain, chest tube, capsule endoscopy} and with \emph{masking, segmentation, background,
region of interest, counterfactual, invariance, concept erasure, group robustness}. We then followed references
backwards and citations forwards from four anchor papers \citep{degrave2021ai,brown2023shortcut,jimenez2023shortcuts,sourget2025mask}.

\paragraph{Inclusion} A study was included if it (i) was published, in final peer-reviewed form, between January 2021
and October 2026; (ii) appeared in a journal or main-track conference of recognised standing; and (iii) addressed one
of RQ1--RQ6. For reading convenience we mark venues as \textbf{A} (Nature, Lancet and Science families, Medical Image
Analysis, Radiology: Artificial Intelligence, JAMIA, Information Fusion; NeurIPS, ICML, ICLR, CVPR, ECCV, AAAI, AISTATS
and MICCAI) or \textbf{B} (other established peer-reviewed journals and conferences, e.g.\ European Radiology, European
Journal of Cancer, JMIR, Computers in Biology and Medicine, Machine Learning, Patterns, MIDL, ISBI, MLSys). The tiers
are an aid for the reader, not a formal ranking.

\paragraph{Exclusion} Preprints without a peer-reviewed version, workshop papers and papers in journals without a
consistent standing in the field were excluded even when relevant. They are listed, uncited, in
Appendix~\ref{app:excluded} so that the reader can see what was left out and why. Studies of segmentation accuracy
alone, and fairness studies without a bearing on the mechanism of shortcut use, were excluded. A small number of works
published before 2021 are cited as background; they are marked with \bg\ and are not counted.

\paragraph{Verification} For every included study the title, authors, venue, volume, pages or article number and DOI
were checked against the publisher's or proceedings' own page. Each summary below states only what the paper's
abstract or main results report; where we draw an inference of our own, the text says so.

\paragraph{Limitations of the method} This is a structured narrative review by one reviewer, not a PRISMA systematic
review: screening counts were not recorded, only English-language work was searched, and coverage of 2026 is
necessarily incomplete because indexing lags publication. Statements of the form ``we found no study'' mean no study
within this search.

%======================================================================================================================
\section{Evidence that medical classifiers learn shortcuts (RQ1)}\label{sec:evidence}
\subsection{Confounded acquisition and dataset assembly}
The early warning came from pneumonia detection, where models identified the hospital of origin from radiographs
\citep{zech2018variable}\bg, and from \emph{hidden stratification}, where overall metrics conceal failures on
clinically important subsets \citep{oakden2020hidden}\bg. The COVID-19 literature then made the problem unmistakable.
\citet{degrave2021ai} showed with saliency maps and generative counterfactuals that radiographic COVID-19 detectors
relied on laterality markers, patient positioning and other confounds of the source datasets rather than on lung
pathology, and lost much of their performance on external data. A systematic review of 62 COVID-19 imaging models found
none suitable for clinical use, citing bias in small datasets and the variability of large, internationally sourced
ones \citep{roberts2021pitfalls}.
\citet{maguolo2021critic} classified COVID-19 radiographs whose lung fields had been blacked out, showing that the
datasets carried label information outside the anatomy the task is about.

\subsection{Demographic attributes as encoded shortcuts}
Deep models predict self-reported race from radiographs, even from degraded images \citep{gichoya2022race}, and
chest-radiograph classifiers underdiagnose under-served groups, most strongly at intersections such as Hispanic
female patients \citep{seyyedkalantari2021underdiagnosis}. Disease models encode protected characteristics in their
features \citep{glocker2023encoding}, and a chest-radiograph foundation model showed sex- and race-related biases
\citep{glocker2023foundation}; expert-level vision--language foundation models underdiagnose marginalised groups more
than radiologists do \citep{yang2025demographic}. \citet{yang2024limits} found that models which encode demographics
more strongly show larger fairness gaps, and that debiasing which is optimal in distribution does not transfer to new
sites. Two methods test reliance rather than encoding: \emph{shortcut testing} varies how strongly a sensitive attribute
is encoded and measures the effect on fairness \citep{brown2023shortcut}, and \citet{ongly2024shortcut} estimate
external performance without external data from a model's dependence on shortcuts. \citet{jones2024causal} separate
three causal mechanisms of dataset bias (prevalence, presentation and annotation) that look alike in the data but call
for different mitigations, and \citet{banerjee2023shortcuts} review causes, evaluation and mitigation of shortcuts for
radiologists. For foundation-model features, pooling data across institutions can hide shortcut-like behaviour that
per-site evaluation reveals \citep{germani2026aggregation}.

\subsection{Controlled, synthetic bias}
Several studies inject a known bias to study its effect. \citet{stanley2024simba} generate synthetic brain images with
known disease and bias effects, compare three mitigations (reweighing worked best) and show how explanations reveal the
bias; a follow-up locates where in a convolutional network the bias is encoded and finds that strong encoding does not
always mean the model uses it \citep{stanley2025where}. \citet{sun2023right} contaminate chest radiographs with a tag,
hyperintensities or an obstruction and find that explanation methods differ strongly in whether they reveal the
spurious signal. \citet{boland2024shortcuts} study shortcuts such as text labels added to images and locate where in the network the
shortcut is used.
\takeaway{The traps of the thesis belong to this controlled tradition, but every synthetic-bias design above fixes where
the bias is placed (or makes it global, such as an intensity shift); none varies the artifact's position relative to the
ROI, which is the variable the thesis manipulates.}

%======================================================================================================================
\section{Artifact shortcuts by modality (RQ2)}\label{sec:modality}
\subsection{Dermoscopy}
Dermoscopy has the richest artifact literature. Ruler scale bars changed the melanoma predictions of a market-approved
classifier \citep{winkler2021scalebars}, as surgical skin markings had earlier \citep{winkler2019markings}\bg.
Explanation-based audits found colour calibration patches, rulers and other artifacts used by ISIC models
\citep{anders2022clever,pahde2023r2r}, and the approach was later extended to semi-automated bias annotation across four
medical tasks \citep{pahde2025ensuring}. Hair occluding the lesion degrades classifiers \citep{wang2024hair}, and adding
synthetic dark-corner artifacts during training helps more than removing real ones \citep{pewton2024dca}. The ISIC
collections contain duplicate images across releases, which can leak between training and test sets
\citep{cassidy2022isic}; the ISIC 2020 set
was released with patient-level metadata \citep{rotemberg2021isic2020}. Performance falls on darker skin and uncommon
diseases \citep{daneshjou2022disparities}. Lesion segmentation, the source of dermoscopy ROIs, is reviewed by
\citet{mirikharaji2023survey}.

\emph{Location.} Calibration patches and dark corners lie mostly outside the lesion; rulers usually lie near the image
edge and markings around the lesion, sometimes overlapping its border; hair crosses it. The published studies report the effect of an artifact's presence, not of where it lies.

\subsection{Ultrasound}
In fetal ultrasound, calipers and on-screen text act as shortcuts even for segmentation \citep{lin2024segshortcut}.
The public thyroid datasets used in the thesis \citep{gong2021tn3k,hou2024thyus2path} contain sonographer calipers,
mostly on or around the nodule (our observation; the dataset papers do not study them). Within our search we found no peer-reviewed study that quantifies
calipers as a \emph{classification} shortcut in thyroid or ovarian ultrasound, or that relates the caliper's position to
the nodule ROI.

\subsection{Chest radiography}
Pneumothorax classifiers detect the chest tube that treats the condition. \citet{rueckel2021pneumothorax} reduced this
confounding by training with in-image annotations, and \citet{jimenez2023shortcuts} labelled
drains to expose the shortcut on a public dataset. A meta-analysis of deep pneumothorax detection reports performance
comparable to physicians but a high risk of bias in most studies \citep{sugibayashi2023pneumothorax}. Saliency maps
localise chest pathology worse than radiologists \citep{saporta2022saliency}, which limits explanation-based audits.
Diffusion-based editing allows stress tests for acquisition, manifestation and population shift without new data
\citep{perezgarcia2024radedit}. \citet{bassi2024isnet} penalise relevance outside segmentation masks to make chest
radiograph classifiers robust to background bias.

\emph{Location.} A drain projects over the lung field, so for lung-field masks it is an in-ROI artifact. Lung
masking therefore cannot remove it: this is the boundary case discussed in the thesis.

\subsection{Capsule endoscopy and other modalities}
\citet{pahde2025ensuring} include a gastrointestinal abnormality task among their four. We found no peer-reviewed study
of bubbles or debris as a classification shortcut in capsule endoscopy, a modality in which the artifact and the lesion
can share the same pixels.

\begin{table}[t]\centering\small
\caption{Artifacts documented as shortcuts (2021--2026), with their typical position relative to the ROI used for
masking. ``Position'' is our reading of the artifact type, not a quantity reported by the cited studies.}\label{tab:artifacts}
\begin{tabular}{P{2.3cm}P{3.6cm}P{3.3cm}P{5.9cm}}\toprule
Modality & Artifact & Typical position & Studies \\\midrule
Dermoscopy & ruler, scale bar, skin marking & mostly outside; markings can overlap the border & \citet{winkler2021scalebars,anders2022clever,pahde2023r2r} \\
 & hair & crosses the lesion & \citet{wang2024hair} \\
 & dark corner, colour patch & outside & \citet{pewton2024dca,anders2022clever} \\
Ultrasound & calipers, on-screen text & on or around the ROI & \citet{lin2024segshortcut} (segmentation only) \\
Chest radiograph & chest tube or drain & inside the lung field & \citet{rueckel2021pneumothorax,jimenez2023shortcuts} \\
 & laterality marker, text, positioning & outside & \citet{degrave2021ai} \\
 & source-dataset signal & everywhere, incl.\ outside & \citet{maguolo2021critic,sourget2025mask} \\
Capsule endoscopy & bubbles, debris & can overlap the lesion & none found \\\bottomrule
\end{tabular}\end{table}

%======================================================================================================================
\section{Restricting the model to the region of interest (RQ3)}\label{sec:masking}
\subsection{Backgrounds carry signal}
In natural images, ImageNet classifiers recognise objects from backgrounds alone, and more accurate models rely less on
the background \citep{xiao2021noise}; foreground and background sensitivity varies strongly with training and
architecture \citep{moayeri2022comprehensive}. The same holds in medicine. COVID-19 radiographs remain classifiable with
the lungs removed \citep{maguolo2021critic}, and in a systematic masking study on chest radiographs and fundus images,
models trained on full images performed better on images \emph{without} the ROI than on images that contained only the
ROI \citep{sourget2025mask}.

\subsection{ROI restriction as a remedy}
Training on segmented dermoscopy images improved balanced accuracy for classifiers trained on HAM10000 but not for those
trained on ISIC \citep{maron2021confounding}. ISNet uses segmentation masks only at training time, penalising relevance
outside them, and improves robustness to background bias on chest radiographs \citep{bassi2024isnet}; in-image
annotations serve a similar purpose for chest tubes \citep{rueckel2021pneumothorax}. Each of these remedies assumes
that the shortcut lies outside the region that is kept.

\subsection{Removing one shortcut can amplify another}
\citet{li2023whacamole} showed in natural images that when several shortcuts are present, mitigating one makes the
model rely more on the others, and proposed a last-layer ensemble against it. In our reading, ROI masking is a
mitigation of exactly this kind: it removes background shortcuts and, if an artifact remains inside the ROI, the same
dilemma predicts greater reliance on it. Li et~al.\ did not study masking or medical images.
\takeaway{The prior work supports both halves of the problem---backgrounds carry shortcuts, and masking often
helps---but every ROI-restriction study assumes the shortcut is in the background. Within our search, no study varies
the artifact's location relative to the ROI, and none tests whether masking increases reliance on an in-ROI artifact
in medical images.}

%======================================================================================================================
\section{Detecting and auditing shortcuts (RQ4)}\label{sec:detect}
\paragraph{Explanations} Spectral relevance analysis clusters explanation heatmaps to reveal Clever-Hans strategies in
ImageNet and ISIC 2019 \citep{anders2022clever}; Reveal-to-Revise turns this into an iterative cycle of revealing,
modelling, removing and re-evaluating a bias \citep{pahde2023r2r,pahde2025ensuring}. Explanations are not reliably
faithful, however: detection of injected confounders differs widely between methods \citep{sun2023right}, saliency
localises chest pathology worse than radiologists \citep{saporta2022saliency}, and shortcut dependence can be traced to
particular layers \citep{boland2024shortcuts}. Our inference: a heatmap cannot separate an in-ROI artifact from the
pathology it overlaps, because both occupy the same pixels.

\paragraph{Counterfactuals and editing} Diffusion-based counterfactuals add or remove a suspected shortcut and measure
the change in prediction, with a masking scheme that confines the edit to a small region \citep{weng2024fastdime};
RadEdit edits chest radiographs under multiple masks to stress-test classifiers \citep{perezgarcia2024radedit}; and
\citet{degrave2021ai} used generative counterfactuals to confirm the COVID-19 confounds.

\paragraph{Statistical tests} Shortcut testing \citep{brown2023shortcut}, estimation of external performance from
internal data \citep{ongly2024shortcut}, artifact labels on public datasets \citep{jimenez2023shortcuts} and
per-subgroup evaluation \citep{yang2024limits} probe reliance without heatmaps.
\takeaway{The thesis uses the counterfactual idea in its simplest form, pasting a real artifact into the same image
inside or outside the ROI, which yields a same-image test of location rather than of presence.}

%======================================================================================================================
\section{Mitigation (RQ5)}\label{sec:mitigate}
\subsection{Group robustness}
With group labels (artifact $\times$ class), GroupDRO \citep{sagawa2020groupdro}\bg, variance-of-risk penalties
(V-REx; \citealp{krueger2021rex}) and last-layer retraining on a group-balanced set (DFR; \citealp{kirichenko2023dfr})
are standard. \citet{izmailov2022feature} showed that ERM features already contain much of the core information, so
retraining only the last layer can recover robustness, which motivates linear heads on frozen features. Without group
labels, JTT up-weights the errors of a first model \citep{liu2021jtt} and AFR reweights the last layer automatically
\citep{qiu2023afr}. Large benchmarks temper expectations: across 20 algorithms and 12 datasets, methods improved
subgroup robustness only for some types of shift \citep{yang2023change}, and in medical imaging, bias-mitigation
algorithms did not significantly improve fairness over ERM, while model selection mattered more
\citep{zong2023medfair}.

\subsection{Auxiliary labels and invariance}
Counterfactual invariance formalises what a shortcut-free predictor should satisfy; in the anti-causal setting, in
which the label causes the image, it implies that the prediction is independent of the nuisance \emph{given the label}
\citep{veitch2021counterfactual}. \citet{makar2022causally} enforce the conditional independences implied by the causal
graph with auxiliary labels available at training time, and NuRD trains under a distribution in which the
nuisance--label relation is broken \citep{puli2022nurd}. \citet{kaur2023modeling} show that the correct independence
constraint depends on the data-generating process and that no single constraint suits all shifts, consistent with the
mechanism-specific view of \citet{jones2024causal}.

\subsection{Concept erasure}
Linear erasure removes a concept from the features: iterative nullspace projection \citep{ravfogel2020inlp}\bg,
adversarial erasure \citep{ravfogel2022rlace} and the closed-form LEACE \citep{belrose2023leace}. When the concept is
correlated with the label, erasure also removes label signal; SPLINCE preserves covariance with the task label while
removing the concept \citep{holstege2025splince}.

\subsection{Explanation-guided and data-centric methods}
Class-artifact compensation \citep{anders2022clever}, gradient penalties along a concept direction in latent space
\citep{dreyer2024hope} and relevance penalties outside a mask \citep{bassi2024isnet} use explanations to suppress a known
bias. On the data side, in-image annotations \citep{rueckel2021pneumothorax}, artifact augmentation
\citep{pewton2024dca}, diffusion-generated synthetic data, which improved fairness out of distribution in
histopathology, chest radiography and dermatology \citep{ktena2024generative}, and counterfactual removal of shortcuts
\citep{weng2024fastdime} change what the model sees.

\begin{table}[t]\centering\small
\caption{Mitigation families and what they need. ``Thesis arm'' names the method of the thesis that belongs to each
family (ERM: standard training; mask: ROI masking).}\label{tab:mitigation}
\begin{tabular}{P{2.7cm}P{3.3cm}P{2.2cm}P{3.5cm}P{2.8cm}}\toprule
Family & Representative studies & Needs at training & Evidence in medical imaging & Thesis arm \\\midrule
ROI restriction & \citet{maron2021confounding,bassi2024isnet} & ROI masks & mixed; assumes background shortcut & mask \\
Group reweighting / DRO & \citet{krueger2021rex,kirichenko2023dfr,qiu2023afr} & group labels (AFR, JTT: none) & no consistent gain over ERM \citep{zong2023medfair} & group balancing \\
Conditional invariance & \citet{veitch2021counterfactual,makar2022causally,puli2022nurd} & auxiliary (artifact) labels & limited & class-conditional constraint \\
Concept erasure & \citet{belrose2023leace,holstege2025splince} & concept labels or pairs & limited & U-MtE (label-free) \\
Explanation-guided & \citet{anders2022clever,dreyer2024hope,pahde2023r2r} & concept annotation & dermoscopy, chest radiograph & --- \\
Generative / data & \citet{ktena2024generative,weng2024fastdime} & a generative model & histopathology, chest radiograph, dermoscopy & --- \\\bottomrule
\end{tabular}\end{table}
\takeaway{The thesis's recommended remedy, a class-conditional mean constraint on the masked features, is the
first-moment, linear case of the conditional independence of \citet{veitch2021counterfactual} and
\citet{makar2022causally}. What the thesis adds is its derivation from the location law, its use after masking, and a
registered test of whether any remedy \emph{dominates} masking; the benchmark results of \citet{yang2023change} and
\citet{zong2023medfair} make it unsurprising that none does.}

%======================================================================================================================
\section{Theory of feature choice (RQ6, part 1)}\label{sec:theory}
Why does a model pick a shortcut when a stable feature is available? \citet{nagarajan2021understanding} show that even
when the invariant feature fully predicts the label, geometric and statistical skews push ERM toward spurious ones.
Gradient descent with cross-entropy can starve less salient features \citep{pezeshki2021gradient}, and the max-margin
bias of default ERM produces shortcut reliance even when the shortcut adds no information \citep{puli2023dontblame}.
\citet{hermann2024foundations} separate a feature's \emph{predictivity} from its \emph{availability} (how easily it is
extracted) and quantify shortcut bias as reliance on the more available but less predictive feature. Invariant risk
minimisation can fail outside the linear regime or with few environments \citep{rosenfeld2021risks}. For frozen
features, linear probing can generalise better out of distribution than fine-tuning, which distorts good pretrained
features \citep{kumar2022finetuning}.
\takeaway{The thesis's linear-Gaussian account fits this line of work. Masking removes competing evidence outside the
ROI, which in Hermann et al.'s terms changes the relative predictivity of what remains, so reliance on an in-ROI
artifact grows. Unlike the explanatory theories above, the account was used to register quantitative predictions
before the results existed.}

%======================================================================================================================
\section{Evaluation methodology (RQ6, part 2)}\label{sec:eval}
Methodological failures in medical-imaging machine learning, from dataset bias to publication incentives, are reviewed
by \citet{varoquaux2022failures}, and data leakage, including non-independent splits, inflated results in at least 294
papers across 17 fields \citep{kapoor2023leakage}. Variance from data sampling, initialisation and hyperparameters
changes benchmark conclusions, and accounting for several sources at once estimates it better
\citep{bouthillier2021variance}. In MICCAI papers, confidence intervals are rarely reported and their typical width
exceeds the gap between the first- and second-ranked methods \citep{christodoulou2024ci}. Metrics should follow from
the clinical question \citep{maierhein2024metrics}, and reporting should follow CLAIM \citep{tejani2024claim}. Model
selection can matter more than the mitigation algorithm \citep{zong2023medfair,yang2023change}, and subgroup failures
are hidden by overall metrics \citep{yang2024limits}, as hidden stratification had shown
\citep{oakden2020hidden}\bg.
\takeaway{These studies motivate the evaluation choices of the thesis: patient- or lesion-grouped splits; a bootstrap
that resamples training seeds and test images jointly, since a test set shared by several seeds is one sample, not
several; Holm-corrected families; registration before results; and evaluation on shortcut-conflicting cases as well
as on the whole unaltered test set.}

%======================================================================================================================
\section{Gap analysis: where the thesis stands}\label{sec:gap}
Table~\ref{tab:gap} sets each contribution of the thesis against the closest peer-reviewed work found.

\begin{longtable}{P{2.2cm}P{4.5cm}P{3.8cm}P{4.5cm}}
\caption{Gap analysis. ``Open'' means not addressed by any included study.}\label{tab:gap}\\\toprule
Question & Established in 2021--2026 & Open & Thesis contribution \\\midrule\endfirsthead
\toprule Question & Established in 2021--2026 & Open & Thesis contribution \\\midrule\endhead
\bottomrule\endfoot
Where does the shortcut sit? & Artifacts are documented by presence \citep{winkler2021scalebars,rueckel2021pneumothorax,lin2024segshortcut,wang2024hair}; backgrounds carry signal \citep{maguolo2021critic,sourget2025mask,xiao2021noise} & Location relative to the ROI as an experimental variable & Overlap $r$ controlled in sweeps; real traps with in-ROI and out-of-ROI artifacts; transplant of the same real artifact across the ROI boundary of the same images \\
Does ROI masking fix it? & Sometimes \citep{maron2021confounding,bassi2024isnet}; every study assumes a background shortcut & When and why masking fails & Location law: masking removes out-of-ROI shortcuts and leaves in-ROI ones, in four cohorts and three modalities \\
Can a mitigation amplify a shortcut? & Yes, for multiple shortcuts in natural images \citep{li2023whacamole} & Whether masking amplifies an in-ROI artifact in medical images, and by what mechanism & Reliance on the in-ROI artifact grows after masking because competing evidence is removed \\
Can it be predicted? & Explanatory theories of feature choice \citep{nagarajan2021understanding,puli2023dontblame,hermann2024foundations} & Quantitative, prospectively tested predictions & Two-parameter linear-Gaussian account; predictions registered before the results \\
Is harm visible clinically? & Overall metrics hide subgroup failures \citep{yang2024limits}; see also \citet{oakden2020hidden}\bg & Effect of masking on shortcut-conflicting cases of unaltered data & Lower AUROC on conflicting cases of patient-disjoint and new-patient thyroid data; sensitivity loss among malignant nodules with calipers \\
What to use instead? & No method uniformly beats ERM \citep{yang2023change,zong2023medfair}; conditional invariance with auxiliary labels \citep{veitch2021counterfactual,makar2022causally}; erasure \citep{belrose2023leace,holstege2025splince} & Remedies tested against masking under a registered criterion & 42-component dominance criterion, two registered rounds; no remedy dominates; the class-conditional constraint comes closest \\
Which modalities? & Mostly dermoscopy and chest radiography; calipers in fetal-US segmentation \citep{lin2024segshortcut} & Thyroid and ovarian ultrasound calipers, capsule debris as classification shortcuts & Traps in thyroid and ovarian ultrasound and capsule endoscopy, alongside dermoscopy \\
How strong is the evidence? & CIs rarely reported \citep{christodoulou2024ci}; seed and data variance \citep{bouthillier2021variance}; leakage \citep{kapoor2023leakage} & Joint seed $\times$ image uncertainty for shortcut experiments & Crossed bootstrap, Holm families, registration, every interval traced to a result file \\
\end{longtable}

\paragraph{What is not new} That a mask cannot remove what it keeps is self-evident; what was not established is that
masking makes reliance on a kept artifact \emph{grow}, how large this is on real artifacts, that it is caused by
location, and how to predict it. Group balancing, last-layer retraining, IRM and V-REx are established methods used as
baselines. The class-conditional constraint is a special case of known conditional-invariance regularisers
\citep{veitch2021counterfactual,makar2022causally}, and disease-preserving erasure overlaps with SPLINCE
\citep{holstege2025splince}; the thesis presents both as adaptations, not as new methods.

\paragraph{Threats to this positioning} (i) Two workshop papers study masking directly: early versus late background
masking for fine-grained recognition, and artifact-based domain generalisation with test-time masking in dermoscopy
(Appendix~\ref{app:excluded}). Both assume background shortcuts, but a reviewer may know them. (ii) The whac-a-mole
dilemma \citep{li2023whacamole} anticipates the amplification conceptually; the thesis should cite it and state that it
supplies the location-specific mechanism and the medical evidence. (iii) Coverage of 2026 is incomplete; the review
should be repeated before submission.

%======================================================================================================================
\section{Summary of included studies}\label{sec:table}
\begingroup\footnotesize\setlength{\tabcolsep}{3pt}
\begin{longtable}{P{2.5cm}P{2.5cm}P{0.6cm}P{1.6cm}P{4.15cm}P{4.2cm}}
\caption{All 68 included studies (2021--2026), grouped by theme. Tier: A or B as defined in Sec.~\ref{sec:method}.
CXR: chest radiograph; US: ultrasound; NI: natural images; --: not modality-specific.}\label{tab:all}\\\toprule
Study & Venue & Tier & Modality & Finding or contribution & Relevance to the thesis \\\midrule\endfirsthead
\toprule Study & Venue & Tier & Modality & Finding or contribution & Relevance to the thesis \\\midrule\endhead
\bottomrule\endfoot
\multicolumn{6}{l}{\textit{Evidence of shortcut learning and bias}}\\
\citet{degrave2021ai} & Nat.\ Mach.\ Intell. & A & CXR & COVID-19 detectors use dataset confounds & motivating example \\
\citet{roberts2021pitfalls} & Nat.\ Mach.\ Intell. & A & CXR, CT & none of 62 models clinically usable & dataset bias \\
\citet{maguolo2021critic} & Inf.\ Fusion & A & CXR & classification with lungs blacked out & signal outside the ROI \\
\citet{seyyedkalantari2021underdiagnosis} & Nat.\ Med. & A & CXR & underdiagnosis of under-served groups & hidden subgroup harm \\
\citet{gichoya2022race} & Lancet Digit.\ Health & A & multiple & race predictable from images & invisible correlates \\
\citet{glocker2023encoding} & eBioMedicine & A & CXR & disease models encode protected attributes & encoding vs use \\
\citet{glocker2023foundation} & Radiol.\ AI & A & CXR & bias in a foundation model & frozen-feature heads \\
\citet{brown2023shortcut} & Nat.\ Commun. & A & derm., CXR & shortcut testing & test of reliance \\
\citet{banerjee2023shortcuts} & J.\ Am.\ Coll.\ Radiol. & B & radiology & review of causes and mitigation & background \\
\citet{yang2024limits} & Nat.\ Med. & A & multiple & fairness does not transfer to new sites & evaluation on new data \\
\citet{yang2025demographic} & Sci.\ Adv. & A & CXR & VLM foundation models underdiagnose & foundation models \\
\citet{jones2024causal} & Nat.\ Mach.\ Intell. & A & -- & three causal bias mechanisms & choice of remedy \\
\citet{ongly2024shortcut} & npj Digit.\ Med. & A & medical & generalisation estimated without external data & reliance estimation \\
\citet{stanley2024simba} & JAMIA & A & brain MRI & synthetic bias framework; reweighing best & controlled design \\
\citet{stanley2025where} & eBioMedicine & A & brain MRI & where bias is encoded & encoding vs use \\
\citet{sun2023right} & MICCAI & A & CXR & explanations vary in detecting injected confounders & limits of heatmaps \\
\citet{boland2024shortcuts} & MIDL & B & medical & layer-wise shortcut localisation & detection \\
\citet{sourget2025mask} & J.\ Imaging Inform.\ Med. & B & CXR, fundus & ROI-removed images remain classifiable & closest masking study \\
\citet{germani2026aggregation} & Med.\ Image Anal. & A & medical & pooled data hide shortcut-like behaviour & foundation features \\
\citet{ktena2024generative} & Nat.\ Med. & A & histo., CXR, derm. & synthetic data improve OOD fairness & generative remedy \\
\multicolumn{6}{l}{\textit{Dermoscopy}}\\
\citet{winkler2021scalebars} & Eur.\ J.\ Cancer & B & derm. & scale bars alter a market-approved CNN & artifact by presence \\
\citet{maron2021confounding} & JMIR & B & derm. & segmentation helps on HAM, not ISIC & masking evidence (mixed) \\
\citet{daneshjou2022disparities} & Sci.\ Adv. & A & derm.\ (clinical) & lower accuracy on dark skin & distribution shift \\
\citet{cassidy2022isic} & Med.\ Image Anal. & A & derm. & ISIC duplicates and benchmarks & grouped splits \\
\citet{mirikharaji2023survey} & Med.\ Image Anal. & A & derm. & survey of lesion segmentation & source of ROI masks \\
\citet{wang2024hair} & Comput.\ Biol.\ Med. & B & derm. & hair degrades classifiers & in-ROI artifact \\
\citet{pewton2024dca} & Comput.\ Methods Programs Biomed. & B & derm. & adding dark corners beats removal & out-of-ROI artifact \\
\citet{anders2022clever} & Inf.\ Fusion & A & derm., NI & SpRAy and class-artifact compensation & detection and remedy \\
\citet{pahde2023r2r} & MICCAI & A & derm., X-ray & Reveal-to-Revise life cycle & detection and remedy \\
\citet{pahde2025ensuring} & Mach.\ Learn. & B & 4 tasks & semi-automated bias annotation & artifact labels \\
\multicolumn{6}{l}{\textit{Ultrasound and chest radiography}}\\
\citet{lin2024segshortcut} & MICCAI & A & US & calipers and text as segmentation shortcuts & ultrasound calipers \\
\citet{rueckel2021pneumothorax} & Eur.\ Radiol. & B & CXR & annotations reduce chest-tube confounding & in-ROI device \\
\citet{sugibayashi2023pneumothorax} & Eur.\ Respir.\ Rev. & B & CXR & meta-analysis; high risk of bias & drain context \\
\citet{jimenez2023shortcuts} & IEEE ISBI & B & CXR & drain labels expose a shortcut & chest boundary case \\
\citet{saporta2022saliency} & Nat.\ Mach.\ Intell. & A & CXR & saliency localises worse than radiologists & limits of heatmaps \\
\citet{perezgarcia2024radedit} & ECCV & A & CXR & diffusion editing for stress tests & counterfactual testing \\
\citet{bassi2024isnet} & Nat.\ Commun. & A & CXR & relevance penalties outside masks & mask-based remedy \\
\multicolumn{6}{l}{\textit{Backgrounds and masking in natural images}}\\
\citet{xiao2021noise} & ICLR & A & NI & backgrounds alone predict the class & background shortcuts \\
\citet{moayeri2022comprehensive} & CVPR & A & NI & foreground/background sensitivity & masking evaluation \\
\citet{li2023whacamole} & CVPR & A & NI & mitigating one shortcut amplifies others & closest to amplification \\
\multicolumn{6}{l}{\textit{Detection by counterfactuals and explanation-guided training}}\\
\citet{weng2024fastdime} & ECCV & A & NI, medical & diffusion counterfactuals add/remove shortcuts & same-image test \\
\citet{dreyer2024hope} & AAAI & A & derm., NI & latent gradient penalisation & explanation remedy \\
\multicolumn{6}{l}{\textit{Mitigation: group robustness and benchmarks}}\\
\citet{krueger2021rex} & ICML & A & -- & V-REx & baseline \\
\citet{liu2021jtt} & ICML & A & -- & JTT, no group labels & label-free baseline \\
\citet{izmailov2022feature} & NeurIPS & A & -- & ERM features hold core information & frozen features \\
\citet{kirichenko2023dfr} & ICLR & A & -- & last-layer retraining (DFR) & balanced head \\
\citet{qiu2023afr} & ICML & A & -- & automatic feature reweighting & label-free baseline \\
\citet{yang2023change} & ICML & A & mixed & no method best across shifts & no dominance expected \\
\citet{zong2023medfair} & ICLR & A & multiple & mitigation $\approx$ ERM in medicine & no dominance expected \\
\multicolumn{6}{l}{\textit{Mitigation: invariance and erasure}}\\
\citet{veitch2021counterfactual} & NeurIPS & A & -- & counterfactual invariance & basis of the constraint \\
\citet{makar2022causally} & AISTATS & A & -- & auxiliary-label conditional independence & closest remedy \\
\citet{puli2022nurd} & ICLR & A & NI, CXR & nuisance-randomised distillation & alternative remedy \\
\citet{kaur2023modeling} & ICLR & A & -- & constraint depends on the causal graph & why class-conditional \\
\citet{ravfogel2022rlace} & ICML & A & -- & adversarial linear erasure & erasure family \\
\citet{belrose2023leace} & NeurIPS & A & -- & closed-form linear erasure & erasure baseline \\
\citet{holstege2025splince} & NeurIPS & A & -- & erasure preserving task signal & disease-preserving erasure \\
\multicolumn{6}{l}{\textit{Theory}}\\
\citet{nagarajan2021understanding} & ICLR & A & -- & skews drive spurious reliance & theory \\
\citet{rosenfeld2021risks} & ICLR & A & -- & limits of IRM & IRM baseline \\
\citet{pezeshki2021gradient} & NeurIPS & A & -- & gradient starvation & theory \\
\citet{kumar2022finetuning} & ICLR & A & -- & linear probing vs fine-tuning OOD & frozen heads \\
\citet{puli2023dontblame} & NeurIPS & A & -- & max-margin causes shortcuts & theory \\
\citet{hermann2024foundations} & ICLR & A & -- & availability vs predictivity & closest to the account \\
\multicolumn{6}{l}{\textit{Evaluation methodology and reporting}}\\
\citet{varoquaux2022failures} & npj Digit.\ Med. & A & -- & methodological failures & study design \\
\citet{kapoor2023leakage} & Patterns & B & -- & leakage taxonomy & grouped splits \\
\citet{bouthillier2021variance} & MLSys & B & -- & sources of benchmark variance & crossed bootstrap \\
\citet{christodoulou2024ci} & MICCAI & A & -- & CIs rarely reported & intervals everywhere \\
\citet{maierhein2024metrics} & Nat.\ Methods & A & -- & problem-driven metrics & metric choice \\
\citet{tejani2024claim} & Radiol.\ AI & A & -- & CLAIM 2024 checklist & reporting \\
\end{longtable}
\endgroup

%======================================================================================================================
\appendix
\section{Identified but excluded by the venue criterion}\label{app:excluded}
These items are relevant but do not meet the venue criterion of Sec.~\ref{sec:method}: preprints, workshop papers or
journals outside it. They are listed for transparency and are not counted.
\begin{itemize}[itemsep=1pt]
\item Aniraj, Dantas, Ienco, Marcos (2023), \emph{Masking strategies for background bias removal in computer vision
models}, ICCV Workshops (OODCV): early (image-level) masking gave the best out-of-distribution accuracy in fine-grained
recognition. The closest general-vision masking study; workshop paper.
\item Bissoto, Barata, Valle, Avila (2022), \emph{Artifact-based domain generalization of skin lesion models}, ECCV
Workshops (ISIC); and Bissoto et al.\ (2023), \emph{Test-time selection for robust skin lesion analysis}, MICCAI
Workshops (ISIC). Workshop papers.
\item Boland et al.\ (2024), \emph{All you need is a guiding hand: mitigating shortcut bias in deep learning models for
medical imaging}, MICCAI FAIMI workshop (LNCS 15198).
\item Nauta, Walsh, Dubowski, Seifert (2022), \emph{Uncovering and correcting shortcut learning in machine learning
models for skin cancer diagnosis}, \emph{Diagnostics} (MDPI).
\item Jin, Gerych, Ghassemi (2024), \emph{MaskMedPaint}, arXiv:2411.10686 (preprint).
\item Pedersen, Sydendal, Cheplygina, Sourget (2026), real-world chest X-ray shortcuts in MedCLIP, arXiv:2608.12086
(preprint).
\item Kina, Petersen (2026), \emph{Prevalence calibration as shortcut mitigation}, arXiv:2609.07922 (preprint).
\item Chuang et al.\ (2023), \emph{Debiasing vision-language models via biased prompts}, arXiv:2302.00070 (preprint).
\item Arjovsky, Bottou, Gulrajani, Lopez-Paz (2019), \emph{Invariant risk minimization}, arXiv:1907.02893: a preprint and
before 2021, but the source of the IRMv1 baseline used in the thesis.
\end{itemize}

\section{Datasets and encoders cited for provenance}\label{app:provenance}
Cited for the data and models the thesis uses, not as evidence: ISIC 2020 \citep{rotemberg2021isic2020}, TN3K
\citep{gong2021tn3k}, the thyroid pathology dataset \citep{hou2024thyus2path}, DINOv2 \citep{oquab2024dinov2},
RAD-DINO \citep{perezgarcia2025raddino} and PanDerm \citep{yan2025panderm}.

\bibliographystyle{plainnat}
\bibliography{refs}
\end{document}
